% Part: normal-modal-logic
% Chapter: axioms-systems
% Section: derived-rules

\documentclass[../../../include/open-logic-section]{subfiles}

\begin{document}

\olfileid{nml}{prf}{der}

\olsection{导出规则}

找出并写出推导过程显然是困难、繁琐且重复的。例如，我们常常想从 $!A \lif !B$ 推出 $\Box !A \lif \Box !B$，即应用规则~\RK。这需要先应用 \Nec，然后写出~\Ax{K} 的适当实例，再应用~\MP。从 $!A \lif !B$ 和 $!B \lif !C$ 推出 $!A \lif !C$，则需要写出（冗长的）重言式实例
\[
(!A \lif !B) \lif ((!B \lif !C) \lif (!A \lif !C))
\]
并应用两次 \MP。我们常常想用一个已知与之等值的公式替换某个子公式，例如，\iftag{prvDiamond}{用 $\lnot\Box\lnot !A$ 替换 $\Diamond !A$}{用 $\lnot\Diamond\lnot !A$ 替换 $\Box !A$}，或者用 $!A$ 替换 $\lnot\lnot !A$。因此，与其写出实际的推导，不如只注明这些中间步骤为何是可推导的更为方便。为此，让我们总结一些关于可推导性的事实。

\begin{prop}
  如果 $\Log{K} \Proves !A_1$, \dots, $\Log{K} \Proves !A_n$，并且 $!B$ 依命题逻辑从 $!A_1$, \dots, $!A_n$ 推出，那么 $\Log{K} \Proves !B$。
\end{prop}

\begin{proof}
  如果 $!B$ 依命题逻辑从 $!A_1$, \dots,~$!A_n$ 推出，那么
  \[
  !A_1 \lif (!A_2 \lif \cdots (!A_n \lif !B)\dots)
  \]
  是一个重言式实例。应用 $n$ 次 \MP 便得到~$!B$ 的一个推导。
\end{proof}

我们将用~\PL 标明此处使用了该命题。

\begin{prop}
  如果 $\Log{K} \Proves !A_1 \lif (!A_2 \lif \cdots (!A_{n-1} \lif
  !A_n)\dots)$ 那么 $\Log{K} \Proves \Box !A_1 \lif (\Box !A_2 \lif
  \cdots (\Box !A_{n-1} \lif \Box !A_n)\dots)$。
\end{prop}

\begin{proof}
  与 \olref[nor]{prop:rk} 的证明类似，通过对 $n$ 进行归纳。
\end{proof}

我们将用~\RK 标明此处使用了该命题。下面举例说明这些结论如何帮助我们更轻松地建立可推导性结论。

\begin{prop}
  $\Log{K} \Proves (\Box!A \land \Box!B) \lif \Box (!A \land !B)$
\end{prop}

\begin{proof}
  \begin{derivation}
    1. & $\Log{K} \Proves !A \lif (!B \lif (!A \land !B))$ & \Taut \\
    2. & $\Log{K} \Proves \Box!A \lif (\Box !B \lif \Box(!A \land !B)))$ & \RK, 1\\
    3. & $\Log{K} \Proves (\Box!A \land \Box!B) \lif \Box(!A \land !B)$ & \PL, 2
  \end{derivation}
\end{proof}

\begin{prop}\ollabel{prop:rewriting}
  如果 $\Log{K} \Proves !A \liff !B$ 并且 $\Log{K} \Proves
  \Subst{!C}{!A}{q}$ 那么 $\Log{K} \Proves \Subst{!C}{B}{q}$
\end{prop}

\begin{proof}
  习题。
\end{proof}

\begin{prob}
  通过对~$!C$ 的复杂度进行归纳来证明 \olref[nml][prf][der]{prop:rewriting}，即证明：如果 $\Log{K} \Proves !A \liff !B$ 那么 $\Log{K} \Proves \Subst{!C}{!A}{q} \liff \Subst{!C}{!B}{q}$。
\end{prob}

当我们想把 $\Diamond$ 转换为 $\Box$（或反之），或者消去公式内部的双重否定时，这个命题特别有用。接下来，每当我们把公式~$!C(!A)$ 替换为~$!C(!B)$ 时，将用“$!A$ for $!B$”标明此处使用了 \olref{prop:rewriting}。换言之，“$!A$ for $!B$” 是以下推导的缩写：

\begin{derivation}
    & $\Proves !C(!A)$ \\
    & $\Proves !A \liff !B$\\
    & $\Proves !C(!B)$ & 依据 \olref{prop:rewriting}
  \end{derivation}

例如：

\begin{prop}
  $\Log{K} \Proves \lnot\Box p \lif \Diamond \lnot p$
\end{prop}

\begin{proof}
  \iftag{prvDiamond}{% Diamond is primitive
  \begin{derivation}
    1. & $\Log{K} \Proves \Diamond \lnot p \liff \lnot\Box\lnot\lnot p$ &
    \Dual\\
    2. & $\Log{K} \Proves \lnot \Box \lnot\lnot p \lif \Diamond \lnot p$ & \PL, 1\\
    3. & $\Log{K} \Proves \lnot\Box p \lif \Diamond \lnot p$ & $p$ for $\lnot\lnot p$\\
  \end{derivation}
  }{% Diamond is defined
  \begin{derivation}
    1. & $\Log{K} \Proves \lnot \Box \lnot\lnot p \lif \Diamond \lnot p$ & \Taut\\
    2. & $\Log{K} \Proves \lnot\Box p \lif \Diamond \lnot p$ & $p$ for $\lnot\lnot p$\\
  \end{derivation}
  第~$1$ 行的公式是 $\lnot q \lif \lnot q$ 的一个实例，其中将 $q$ 代以 $\Box\lnot\lnot p$。$\Diamond\lnot p$ 与 $\lnot\Box\lnot\lnot p$ 是同一公式，因为 $\Diamond$ 定义为~$\lnot\Box\lnot$。}
\end{proof}

在上面的推导中，最后一步“$p$ for $\lnot\lnot p$” 是以下推导的缩写：

\begin{derivation}
    & $\Log{K} \Proves \lnot \Box \lnot\lnot p \lif \Diamond \lnot
    p$\\
    & $\Log{K} \Proves \lnot\lnot p \liff p$ & \Taut\\
    & $\Log{K} \Proves \lnot\Box p \lif \Diamond \lnot p$ & 依据 \olref{prop:rewriting}
  \end{derivation}

在这里，\olref{prop:rewriting} 中的 $!C(q)$、$!A$ 和 $!B$ 分别对应 $\lnot \Box q \lif \Diamond \lnot p$、$\lnot\lnot p$ 和~$p$。

当一个公式包含子公式 $\lnot\Diamond !A$ 时，我们可以利用 \olref{prop:rewriting} 将其替换为 $\Box\lnot !A$，因为 $\Log{K} \Proves \lnot\Diamond !A \liff \Box\lnot !A$。我们将直接用“$\Box\lnot$ for $\lnot\Diamond$”标示这类及类似的替换。

下述命题表明，我们可以以图式的方式建立可推导性结论。例如，前一个命题不仅确立了 $\Log{K} \Proves \lnot\Box p \lif
\Diamond \lnot p$，而且对任意的~$!A$ 确立了 $\Log{K} \Proves \lnot\Box !A \lif \Diamond
\lnot !A$。

\begin{prop}
  如果 $!A$ 是 $!B$ 的一个代入实例并且 $\Log{K} \Proves !B$，那么 $\Log{K} \Proves !A$。
\end{prop}

\begin{proof}
  验证对整个推导应用代入后仍能得到正确的推导，这个过程虽然繁琐但属于常规操作（可通过对~$!B$ 的推导长度进行归纳来完成）。具体来说，重言式实例的代入实例仍是重言式实例；\iftag{prvDiamond}{\Dual{} 和~}{}公理~\Ax{K} 实例的代入实例仍是\iftag{prvDiamond}{\Dual{} 和~}{}\Ax{K} 的实例；并且当在前提和结论中将命题变元代以公式时，\MP{} 和~\Nec{} 规则的应用仍然正确。
\end{proof}

\end{document}